\chapter{[标准库] math (标量)数学库}

\section{数学常数}
\begin{lstlisting}[style=demostyle,numbers=none]
import math
math.pi  # 圆周率 π ≈ 3.141592653589793
math.e   # 自然对数底数 e ≈ 2.718281828459045
\end{lstlisting}

\section{三角函数相关}
\subsection{角度转换}
\begin{lstlisting}[style=demostyle,numbers=none]
import math
math.degrees(x)  # 弧度转角度（180°/π×x）
math.radians(x)  # 角度转弧度（π/180°×x）


# 例如：
math.radians(180)  # 180° = π
math.degrees(math.pi)  # π = 180°
\end{lstlisting}

\subsection{基本三角函数}
\begin{lstlisting}[style=demostyle,numbers=none]
import math
math.sin(x)    # 正弦函数（x为弧度）
math.cos(x)    # 余弦函数
math.tan(x)    # 正切函数
math.hypot(x,y) # 直角三角形斜边 √(x²+y²)
\end{lstlisting}

\subsection{反三角函数}
\begin{lstlisting}[style=demostyle,numbers=none]
import math

math.asin(x)   # 反正弦（返回弧度）
math.acos(x)   # 反余弦
math.atan(x)   # 反正切
\end{lstlisting}

\section{指数与对数相关}
\begin{lstlisting}[style=demostyle,numbers=none]
math.exp(x)    # e的x次方（等价于 e^x）
math.log(x)    # 自然对数 ln(x)
math.log1p(x)  # ln(1+x)（更精确的x接近0时的计算）
math.log10(x)  # 常用对数 lg(x)
math.log2(x)   # 二进制对数 log₂(x)
\end{lstlisting}


\section{幂函数与开方}
\begin{lstlisting}[style=demostyle,numbers=none]
math.sqrt(x)   # 平方根 √x
math.pow(x,y)  # x的y次方（等价于x**y）
math.cbrt(x)   # 立方根 ∛x
math.isqrt(x)  # 整数平方根（返回整数部分）
\end{lstlisting}

\section{特殊功能函数}
\subsection{取整相关}
\begin{lstlisting}[style=demostyle,numbers=none]
math.ceil(x)   # 向上取整（返回≥x的最小整数）
math.floor(x)  # 向下取整（返回≤x的最大整数）
math.trunc(x)  # 截断取整（直接去除小数部分）
\end{lstlisting}

\subsection{数值处理}
\begin{lstlisting}[style=demostyle,numbers=none]
math.fabs(x)   # 绝对值（返回浮点数）
math.factorial(x) # 阶乘 x!
math.fsum(iterable) # 高精度迭代求和
math.modf(x)   # 分解为(小数部分, 整数部分)元组
\end{lstlisting}

\subsection{特殊判断}
\begin{lstlisting}[style=demostyle,numbers=none]
math.isinf(x)  # 判断是否为无穷大
math.isnan(x)  # 判断是否为NaN（非数值）
\end{lstlisting}
